\chapter{Bit timing}
\label{ch:tajming}

{\large\itshape How nodes without a shared clock read the same bit at the same time.\par}

\begin{chapterbox}{In this chapter}
\begin{itemize}
  \item Why a CAN bit is divided into segments instead of just being counted.
  \item The sample point, and why it lies late in the bit period.
  \item Synchronisation: hard at SOF, soft at every edge after that.
  \item What happens when two nodes' oscillators do not run at the same speed.
\end{itemize}
\end{chapterbox}

\section{The problem}

Every node has its own clock. No wire carries the beat --- CAN has four signals and none of them is
a clock.

Even so, every node has to agree on where each bit begins and ends, and at bit level at that:
arbitration in \kapref{ch:arbitrering} relies on two nodes reading the \emph{same} bit interval at
the same time and drawing conclusions from what they see. If they are out of phase they compare
different bits, and the result is arbitrary.

The solution has two parts: a shared \emph{definition} of where in the bit period to read, and a
mechanism that pulls drifting nodes back.

\section{The bit period in segments}

A CAN bit is not an undivided interval but is split into four segments, measured in
\textbf{time quanta} --- short, equally long steps that the node counts with its own clock.

\bookfigure{bit_timing}{The four segments of the bit period. The sample point lies between phase
segments~1 and 2, late in the period.}{fig:timing}

\begin{booktable}{@{}llL@{}}
\thead{Segment} & \thead{Length} & \thead{What it is for} \\ \midrule
Sync segment & 1 tq & This is where an edge is expected. \\
Propagation segment & 1--8 tq & Covers the signal's travel time out along the bus and back. \\
Phase segment 1 & 1--8 tq & Before the sample point; can be lengthened on resynchronisation. \\
Phase segment 2 & 1--8 tq & After the sample point; can be shortened. \\
\end{booktable}

The \textbf{sample point} lies between phase segments 1 and 2. That is where, and only where, the
bit is read.

\subsection{Why late in the period}

The sample point is typically placed 70--80~\% into the bit period. The reason is the propagation
segment, and it is the same reason that makes the length of the bus limit the speed.

During arbitration, a transmitting node must have time to read back what \emph{the farthest} node
drives, before the bit ends. The signal has to make it out and back. The propagation segment is
that time, and the sample point has to come after it --- otherwise the node reads the bus before
the other node's bit has arrived, and arbitration goes wrong.

\begin{aside}
\note That is why bit rate and cable length are linked (\kapref{ch:buss}). A higher speed gives a
shorter bit period, so a shorter propagation segment in absolute time, so a shorter permissible
cable. The relationship is not electrical but logical: it is about what arbitration requires.
\end{aside}

\section{Synchronisation}

\subsection{Hard synchronisation}

At SOF --- the only edge that every node sees at the same time after the bus has been idle --- every
receiver performs a \textbf{hard synchronisation}: it resets its bit counter, wherever it was.

That is once per frame, and it is the only time the nodes may set the clock entirely freely.

\subsection{Resynchronisation}

During the frame that is not enough. Two oscillators that differ by a fraction of a percent drift
apart slowly, and over a frame of 130 bits it becomes measurable.

That is why every receiver resynchronises on \emph{every} recessive-to-dominant edge in the stream.
If the edge comes earlier than expected, phase segment 2 is shortened; if it comes later, phase
segment 1 is lengthened. The adjustment is limited to a number of time quanta called the
\emph{synchronization jump width}.

\textbf{And now bit stuffing falls into place.} The rule of five in \kapref{ch:stuffing} guarantees
that no more than five bit periods ever pass without an edge, which puts an upper bound on how far
two nodes can drift apart before they are pulled back.

\begin{aside}
So the two mechanisms are not independent: bit stuffing exists \emph{for} resynchronisation's sake.
A protocol without stuffing would need either a clock line or considerably more accurate
oscillators.
\end{aside}

\section{Working out the timing}

For a design, two numbers are derived from the system clock and the desired bit rate:

\[ \text{ticks per bit} = \frac{f_{\text{system}}}{\text{bit rate}} \]

With a 50~MHz system clock and 1~Mbit/s that is 50 clock cycles per CAN bit. The sample point at
70~\% is then tick 35.

\begin{cppcode}
// Derived, never written down as a literal: change the bit rate and both follow.
constexpr unsigned TicksPerBit{SystemClockHz / BitRateHz};
constexpr unsigned SampleTick{TicksPerBit * 7U / 10U};
\end{cppcode}

That they are \emph{derived} rather than written in is the point. A design where \code{50} and
\code{35} are written in, each in its own place, is a design where one of them gets changed without
the other.

\section{What happens when it goes wrong}

Two faults are worth recognising, because they look different:

\textbf{Wrong bit rate.} The nodes are out of phase from the first bit. The symptom is that nothing
works at all --- no frames are received, and the transmitting node reports errors straight away.
It is the kinder fault, because it shows immediately.

\textbf{Wrong sample point.} The nodes agree on most bits but not all. The symptom is sporadic CRC
and stuff errors that get worse with a longer cable and higher load. This is the nasty fault,
because it looks like a bad cable.

\section*{Exercises}

\exercise{Work out the timing}{Calculation}
\label{ex:5:1}
A design has a 50~MHz system clock.
\begin{enumerate}[a)]
  \item How many clock cycles are there in a CAN bit at 1~Mbit/s?
  \item At 500~kbit/s?
  \item Where is the sample point, counted in clock cycles, at 1~Mbit/s and 70~\%?
\end{enumerate}

\exercise{The sample point}{Reasoning}
\label{ex:5:2}
\begin{enumerate}[a)]
  \item Why does the sample point lie late in the bit period instead of in the middle?
  \item What would go wrong during arbitration if it were at 20~\%?
  \item Which part of the bit period would have to grow if the cable became twice as long?
\end{enumerate}

\exercise{Stuffing and sync}{Reasoning}
\label{ex:5:3}
\begin{enumerate}[a)]
  \item At most how many bit periods can pass without an edge in a CAN frame?
  \item Which mechanism guarantees that?
  \item What would happen to resynchronisation if bit stuffing were removed?
\end{enumerate}

\exercise{Two symptoms}{Reasoning}
\label{ex:5:4}
Two nodes are supposed to talk to each other.
\begin{enumerate}[a)]
  \item Node A is set to 500~kbit/s and node B to 1~Mbit/s. What do you see?
  \item Both run at 1~Mbit/s, but A samples at 50~\% and B at 80~\%. What do you see, and why is it
    harder to debug?
\end{enumerate}
